In~\cite{knuth1993johann} Donald Knuth provides a remarkable formula for
sums of odd powers,
in terms of central factorial numbers of the second kind $T(n,k)$
\begin{proposition}[Knuth's odd powers sum]
  \label{prop:knuth-odd-powers-sum}
  For non-negative integers $n,m$
  \begin{align*}
    \KnuthRFoldSum{1}{n}{2m-1} = \sum_{k=1}^{m} (2k-1)! \binom{n+k}{2k} T(2m,2k),
  \end{align*}
  where $T(2m,2k)$ are central factorial numbers of the second kind.
\end{proposition}
It is quite interesting to notice that Proposition~\eqref{prop:knuth-odd-powers-sum}
originates from centered
Newton's formula for powers~\eqref{lem:central-newtons-formula-for-powers}
evaluated at zero.
By Newton's formula, for $m\geq 1$, we have
\begin{align*}
  n^{2m} = \sum_{k=1}^{2m} \frac{\centralFactorial{n}{k}}{k!} \cdot \delta^{k} 0^{2m}
\end{align*}
The iteration for $k=0$ can be omitted because $\delta^{0} 0^{2m} = 0$.
Now, we observe that central difference $\delta^k 0^m$ satisfies the parity of
arguments $m$ and $k$, such that
\begin{align*}
  \begin{cases}
    \delta^{k} 0^m \neq 0, \quad &\text{when} \quad m \equiv k \pmod{2}, \\
    \delta^{k} 0^m = 0, \quad &\text{when} \quad m \not\equiv k \pmod{2},
  \end{cases}
\end{align*}
meaning that $\delta^{k} 0^m$ is zero whether $k,m$ are odd and even, and vise versa.
Note that the parity property holds only for central difference of powers
evaluated in zero, by contrast $\delta^{3} 1^{4} = 24  \neq 0$.
Thus, the parity of $\delta^{k} 0^{2m}$ allows us
to omit odd iterations of $k$, because $\delta^{k} 0^{2m}$ is zero for odd $k$
\begin{align*}
  n^{2m} = \sum_{k=1}^{m} \frac{\centralFactorial{n}{2k}}{(2k)!} \cdot \delta^{2k} 0^{2m}.
\end{align*}
By definition of central factorial numbers of the second kind, yields
\begin{align*}
  n^{2m} = \sum_{k=1}^{m} \centralFactorial{n}{2k} T(2m,2k).
\end{align*}
Thus,
\begin{align*}
  n^{2m} = \sum_{k=1}^{m} (2k)! \frac{\centralFactorial{n}{2k}}{(2k)!} T(2m,2k)
\end{align*}
By simplifying $\frac{\centralFactorial{n}{2k}}{(2k)!}$
using Lemma~\eqref{lem:central-factorials-binomial-form}, we get
\begin{align*}
  n^{2m}
  &= \sum_{k=1}^{m} (2k)!  \frac{n}{2k} \binom{n+k-1}{2k-1} T(2m,2k) \\
  &= \sum_{k=1}^{m} (2k-1)!  n \binom{n+k-1}{2k-1} T(2m,2k).
\end{align*}
By factoring out $n$, we obtain
\begin{align*}
  n^{2m-1}
  &= \sum_{k=1}^{m} (2k-1)! \binom{n+k-1}{2k-1} T(2m,2k).
\end{align*}
By hockey stick identity $\sum_{k=0}^{n} \binom{k}{j} = \binom{n+1}{j+1}$
formula for sums of powers yields
\begin{align*}
  \KnuthRFoldSum{1}{n}{2m-1} = \sum_{k=1}^{m} (2k-1)! \binom{n+k}{2k} T(2m,2k).
\end{align*}
The formula above matches the Proposition~\eqref{prop:knuth-odd-powers-sum} precisely.
Clearly, the formula for sums of odd powers can be generalized to multifold case,
by applying hockey-stick identity $\sum_{k=0}^{n} \binom{k}{j} = \binom{n+1}{j+1}$
repeatedly
\begin{proposition}[Knuth's odd powers sum multifold]
  \label{prop:knuth-odd-powers-sum-multifold}
  For non-negative integers $r,n,m$
  \begin{align*}
    \KnuthRFoldSum{r+1}{n}{2m-1}
    = \sum_{k=1}^{m} (2k-1)! \binom{n+k+r}{2k+r} T(2m,2k),
  \end{align*}
  where $T(2m,2k)$ are central factorial numbers of the second kind.
\end{proposition}
For example,
\begin{align*}
  \KnuthRFoldSum{1}{n}{1} &= \tbinom{n+1}{2} 1, \\
  \KnuthRFoldSum{1}{n}{3} &= \tbinom{n+2}{4} 6    + \tbinom{n+1}{2} 1, \\
  \KnuthRFoldSum{1}{n}{5} &= \tbinom{n+3}{6} 120  + \tbinom{n+2}{4} 30   + \tbinom{n+1}{2}, \\
  \KnuthRFoldSum{1}{n}{7} &= \tbinom{n+4}{8} 5040 + \tbinom{n+3}{6} 1680 + \tbinom{n+2}{4} 126 + \tbinom{n+1}{2} 1.
\end{align*}
While multifold sums of odd powers are,
\begin{align*}
  \KnuthRFoldSum{r}{n}{1} &= \tbinom{n+0+r}{1+r} 1, \\
  \KnuthRFoldSum{r}{n}{3} &= \tbinom{n+1+r}{3+r} 6    + \tbinom{n+0+r}{1+r} 1, \\
  \KnuthRFoldSum{r}{n}{5} &= \tbinom{n+2+r}{5+r} 120  + \tbinom{n+1+r}{3+r} 30   + \tbinom{n+0+r}{1+r} 1, \\
  \KnuthRFoldSum{r}{n}{7} &= \tbinom{n+3+r}{7+r} 5040 + \tbinom{n+3+r}{5+r} 1680 + \tbinom{n+1+r}{3+r} 126 + \tbinom{n+0+r}{1+r} 1.
\end{align*}
The coefficients $1, 6, 1, 120, 30, 1,\ldots$
is the sequence \href{https://oeis.org/A303675}{\texttt{A303675}}
in the OEIS~\cite{sloane2003line}.
